% Include this file after the dissertation's \appendix command.
% Required packages: booktabs, listings, tabularx, and xurl.
\providecommand{\tightlist}{%
  \setlength{\itemsep}{0pt}\setlength{\parskip}{0pt}}
\lstdefinestyle{sagepython}{
  language=Python,
  basicstyle=\ttfamily\scriptsize,
  breaklines=true,
  breakatwhitespace=false,
  columns=fullflexible,
  keepspaces=true,
  frame=single,
  framerule=0.4pt,
  showstringspaces=false,
  tabsize=4
}

\chapter{Native-Action Tools Generated by SAGE}
\label{app:native-action-generated-tools}

\section{Purpose and Scope}\label{purpose-and-scope}

This appendix documents the complete set of generated tools from the
validated 1,032-task SAGE run that were authorized to execute a
preserved native ToolSandbox action. The appendix reports the accepted
implementation of each tool exactly as stored in the generated-tool
registry, explains its operational logic, and provides the validation
and contribution evidence needed to interpret its role in the
experiment.

The run began with an empty generated-tool registry. SAGE used
\texttt{gpt-4o-mini} for the actor, automated participant, tool
generation, and model-authored repair. Scenario-name birth and routing
were disabled, synthetic bridge completions were disabled, and generated
tools were selected through the ordinary actor tool-use interface.
Consequently, the tools below were not preinstalled and were not invoked
through a diagnostic force-call mechanism.

The evidence source is the completed run:

\path{outputs/production_cleanup_20260802/full_optional_cleanup_equivalence_20260802_172055/online_build_full_20260802_172101}

The generated-tool registry is:

\path{artifacts/production_cleanup_20260802/full_optional_cleanup_equivalence_20260802_172055_registry/registry_manifest.json}

The lifecycle attribution record is:

\path{artifacts/production_cleanup_20260802/full_optional_cleanup_equivalence_20260802_172055_registry/tool_lifecycle.json}

\section{Definition of a Native-Action Generated
Tool}\label{definition-of-a-native-action-generated-tool}

For this appendix, a generated tool is classified as a
\emph{native-action tool} only when its accepted immutable specification
contains \texttt{native\_action\_delegation: true} and its
implementation directly invokes at least one approved native environment
action. This definition excludes tools that merely calculate values,
select records, or return the name and arguments of a proposed
downstream action.

Native-action generation does not replace the
environment\textquotesingle s implementation of state mutation. The
generated code performs a bounded deterministic operation, such as
normalizing visible inputs or selecting one unambiguous record, and then
delegates the mutation to the preserved ToolSandbox function. The native
function remains responsible for enforcing environment behavior,
recording the tool call, and modifying state.

Admission required the following properties:

\begin{enumerate}
\def\labelenumi{\arabic{enumi}.}
\tightlist
\item
  Every positive validation path invoked exactly one approved native
  action.
\item
  Negative or ambiguous cases invoked no native action and returned an
  explicit abstention.
\item
  The candidate passed source, held-out, and negative applicability
  examples.
\item
  The candidate passed source, schema, output-shape, runtime-smoke, and
  side-effect-preservation checks.
\item
  Native action selection was determined from visible tool inputs rather
  than benchmark identifiers, hidden expected answers, or scenario
  names.
\end{enumerate}

The six tools below satisfied this structural definition. Tool names
beginning with \texttt{prepare\_} or \texttt{plan\_} are included when
their stored implementation actually invokes a native action;
classification is based on executable behavior, not naming convention.

\section{Summary of Generated Native-Action
Tools}\label{summary-of-generated-native-action-tools}

\begin{table}[htbp]
\centering
\scriptsize
\setlength{\tabcolsep}{3pt}
\renewcommand{\arraystretch}{1.15}
\begin{tabularx}{\textwidth}{p{3.0cm} p{2.35cm} c c c X}
\toprule
\textbf{Generated tool} & \textbf{Native action} & \textbf{S/H/N} &
\textbf{Called rows} & \textbf{Outcome delta} & \textbf{Decision} \\
\midrule
\path{prepare_add_contact_args} & \path{add_contact} & 1/1/1 & 7 & -0.2000 & Needs route repair \\
\path{prepare_reminder_creation_args} & \path{add_reminder} & 2/2/4 & 96 & +0.8819 & Retain \\
\path{apply_single_device_state_action} & One device setter & 1/3/2 & 23 & -0.0037 & Retain with route repair \\
\path{select_message_counterparty_for_contact_update} & \path{modify_contact} & 1/2/4 & 29 & +0.5893 & Retain \\
\path{select_action_target_by_recency} & \path{modify_reminder} or \path{remove_reminder} & 1/2/3 & 30 & +0.6818 & Retain with route repair \\
\path{plan_contact_update_from_id} & \path{modify_contact} & 1/1/1 & 7 & -0.0476 & Needs route repair \\
\bottomrule
\end{tabularx}
\caption{Accepted native-action generated tools and called-subset evidence. S/H/N denotes source, held-out, and negative validation cases. Outcome delta is SAGE minus matched control on naturally called rows.}
\label{tab:native-action-generated-tools}
\end{table}

The \emph{naturally called task rows} column counts matched evaluation
rows in which the generated tool appeared in the actor\textquotesingle s
recorded generated-tool calls. It is not the registry\textquotesingle s
internal \texttt{reuse\_count}, which also reflects broader runtime
routing and reuse bookkeeping. The paired outcome delta is the SAGE
outcome minus the matched control outcome for rows where the tool was
called. These values are descriptive attribution evidence; they do not
isolate a causal effect when multiple tools are called in the same task.

All six tools passed their admission validation with no recorded
validation errors and no runtime-smoke failure. None recorded a helper
side-effect incident in the lifecycle audit. Their downstream
contribution was heterogeneous: three produced strong positive
called-subset outcome deltas, one was approximately neutral, and two
were retained as evidence of successful generation but marked for
routing repair because their called subsets were not beneficial in this
run.

\section{Common Execution Contract}\label{common-execution-contract}

Each native-action tool returns a structured record with the following
core fields:

\begin{table}[htbp]
\centering
\footnotesize
\begin{tabularx}{\textwidth}{p{3.0cm} X}
\toprule
\textbf{Field} & \textbf{Meaning} \\
\midrule
\path{status} & \path{success} when one approved native action was executed; \path{abstain} when execution was not justified. \\
\path{confirmation} & Actor-facing summary of the completed operation. \\
\path{abstain_reason} & Explicit reason no native action was executed. \\
\path{native_action} & Name of the preserved native action that was invoked; blank on abstention. \\
\path{native_result} & Unmodified result returned by the preserved native action. \\
\bottomrule
\end{tabularx}
\caption{Common output contract for native-action generated tools.}
\label{tab:native-action-output-contract}
\end{table}

This contract makes the action boundary observable. The generated tool
cannot claim success merely by returning a textual answer: a successful
result must contain the native action name and the native
function\textquotesingle s returned execution result.

\section{Tool 1: Contact Creation}
\label{sec:appendix-prepare-add-contact}

\subsection{Purpose}\label{purpose}

This tool addresses contact-creation failures caused by incomplete or
incorrectly formatted arguments. It requires a visible contact name and
phone number, normalizes common phone-number formatting characters,
conditionally includes a relationship value, and delegates the state
change to \texttt{add\_contact}.

\subsection{Execution Logic}\label{execution-logic}

\begin{enumerate}
\def\labelenumi{\arabic{enumi}.}
\tightlist
\item
  Abstain if either the name or phone number is missing.
\item
  Remove spaces, parentheses, and hyphens from the phone number.
\item
  Construct the native action arguments from the normalized visible
  values.
\item
  Include \texttt{relationship} only when the actor supplies a nonempty
  value.
\item
  Invoke \texttt{add\_contact} exactly once and return its native
  result.
\end{enumerate}

The tool does not search for a contact or infer missing identity
information. Its applicability is limited to requests where the required
values are already visible and unambiguous.

\subsection{Accepted Generated Code}\label{accepted-generated-code}

\begin{lstlisting}[style=sagepython,caption={Accepted contact-creation native-action tool.},label={lst:sage-prepare-add-contact}]
def prepare_add_contact_args(user_request: str, name: str, phone_number: str, relationship: str):
    if not name or not phone_number:
        return {'status': 'abstain', 'confirmation': '', 'abstain_reason': 'Name and phone number are required.', 'native_action': '', 'native_result': None}
    normalized_phone_number = phone_number.replace(' ', '').replace('(', '').replace(')', '').replace('-', '')
    downstream_tool_kwargs = {'name': name, 'phone_number': normalized_phone_number}
    if relationship:
        downstream_tool_kwargs['relationship'] = relationship
    native_result = add_contact(**downstream_tool_kwargs)
    return {'status': 'success', 'confirmation': f'{name} added with phone number {normalized_phone_number}.', 'abstain_reason': '', 'native_action': 'add_contact', 'native_result': native_result}
\end{lstlisting}

\subsection{Validation and
Interpretation}\label{validation-and-interpretation}

The accepted implementation passed one source example, one held-out
example, one negative applicability example, and runtime smoke
validation. Its SHA-256 code hash is
\path{7e5e4bb4b21c7fe745bc679fc0aff95d75bc4b1d27e811162aa1b517bb6617f9}.
It was naturally called in seven matched task rows with no
generated-tool runtime failure or side-effect incident. Its mean paired
outcome delta on those rows was -0.2000, and lifecycle reflection
classified it as requiring routing repair. The negative contribution
does not invalidate the generation mechanism, but it prevents this tool
from being presented as evidence of positive outcome improvement.

\section{Tool 2: Reminder Creation}
\label{sec:appendix-prepare-reminder-creation}

\subsection{Purpose}\label{purpose-1}

This tool converts already-resolved reminder information into one
validated native reminder-creation action. It enforces required content
and timestamp fields, distinguishes optional from required location
information, and avoids creating a reminder when location resolution
remains incomplete.

\subsection{Execution Logic}\label{execution-logic-1}

\begin{enumerate}
\def\labelenumi{\arabic{enumi}.}
\tightlist
\item
  Abstain when reminder content is empty or the timestamp is
  nonpositive.
\item
  Abstain when location is required but unavailable.
\item
  Omit optional coordinates when a requested location lookup failed.
\item
  Abstain when a location lookup is still pending.
\item
  Invoke \texttt{add\_reminder} exactly once with the resolved content,
  timestamp, and permissible coordinates.
\end{enumerate}

The tool does not calculate relative dates or fabricate coordinates.
Those values must be obtained through visible upstream reasoning or
native tool results before invocation.

\subsection{Accepted Generated Code}\label{accepted-generated-code-1}

\begin{lstlisting}[style=sagepython,caption={Accepted reminder-creation native-action tool.},label={lst:sage-prepare-reminder-creation}]
def prepare_reminder_creation_args(content: str, reminder_timestamp: float, location_requested: bool, location_required: bool, location_available: bool, latitude: float, longitude: float, location_lookup_failed: bool):
    if not content:
        return {'status': 'abstain', 'confirmation': '', 'abstain_reason': 'Content is missing.', 'native_action': '', 'native_result': None}
    if reminder_timestamp <= 0:
        return {'status': 'abstain', 'confirmation': '', 'abstain_reason': 'Invalid reminder timestamp.', 'native_action': '', 'native_result': None}
    if location_required and (not location_available):
        return {'status': 'abstain', 'confirmation': '', 'abstain_reason': 'Required location is unavailable.', 'native_action': '', 'native_result': None}
    if location_requested and location_lookup_failed:
        latitude = None
        longitude = None
    elif location_requested and (not location_lookup_failed) and (not location_available):
        return {'status': 'abstain', 'confirmation': '', 'abstain_reason': 'Location lookup is still pending.', 'native_action': '', 'native_result': None}
    native_result = add_reminder(content=content, reminder_timestamp=reminder_timestamp, latitude=latitude, longitude=longitude)
    return {'status': 'success', 'confirmation': 'Reminder created successfully.', 'abstain_reason': '', 'native_action': 'add_reminder', 'native_result': native_result}
\end{lstlisting}

\subsection{Validation and
Interpretation}\label{validation-and-interpretation-1}

The implementation passed two source examples, two held-out examples,
four negative applicability examples, and runtime smoke validation. Its
SHA-256 code hash is
\path{9fe20eee4678d2c81081ef0adf5682f3c58a8e0354811947a63adb5aff7dfea6}.
It was naturally called in 96 matched task rows, recorded no runtime or
side-effect incident, and produced a mean paired outcome delta of
+0.8819 on called rows. Ninety called rows were classified as helpful
and none as harmful. Lifecycle reflection therefore retained the tool.

\section{Tool 3: Single Device-State Action}
\label{sec:appendix-apply-device-state}

\subsection{Purpose}\label{purpose-2}

This tool maps one explicit device service and one explicit Boolean
target state to exactly one approved native setter. It addresses
failures in which the desired state is understood but the final native
state-changing call is not executed correctly.

\subsection{Execution Logic}\label{execution-logic-2}

\begin{enumerate}
\def\labelenumi{\arabic{enumi}.}
\tightlist
\item
  Match \texttt{target\_service} against the finite supported set:
  cellular, location, low-battery mode, or Wi-Fi.
\item
  Invoke the corresponding preserved setter with \texttt{desired\_on}.
\item
  Return the native result and a service-specific confirmation.
\item
  Abstain without a native call when the service value is unsupported.
\end{enumerate}

The tool intentionally performs only one action. It does not execute
recovery sequences, infer additional service dependencies, or change
more than one setting per invocation.

\subsection{Accepted Generated Code}\label{accepted-generated-code-2}

\begin{lstlisting}[style=sagepython,caption={Accepted single-device-state native-action tool.},label={lst:sage-apply-device-state}]
def apply_single_device_state_action(target_service: str, desired_on: bool):
    if target_service == 'cellular':
        native_result = set_cellular_service_status(on=desired_on)
        return {'status': 'success', 'confirmation': 'Cellular service has been updated.', 'abstain_reason': '', 'native_action': 'set_cellular_service_status', 'native_result': native_result}
    elif target_service == 'location':
        native_result = set_location_service_status(on=desired_on)
        return {'status': 'success', 'confirmation': 'Location service has been updated.', 'abstain_reason': '', 'native_action': 'set_location_service_status', 'native_result': native_result}
    elif target_service == 'low_battery_mode':
        native_result = set_low_battery_mode_status(on=desired_on)
        return {'status': 'success', 'confirmation': 'Low battery mode has been updated.', 'abstain_reason': '', 'native_action': 'set_low_battery_mode_status', 'native_result': native_result}
    elif target_service == 'wifi':
        native_result = set_wifi_status(on=desired_on)
        return {'status': 'success', 'confirmation': 'WiFi service has been updated.', 'abstain_reason': '', 'native_action': 'set_wifi_status', 'native_result': native_result}
    else:
        return {'status': 'abstain', 'confirmation': '', 'abstain_reason': 'The target service is unknown or invalid.', 'native_action': '', 'native_result': None}
\end{lstlisting}

\subsection{Validation and
Interpretation}\label{validation-and-interpretation-2}

The implementation passed one source example, three held-out examples,
two negative applicability examples, and runtime smoke validation. Its
SHA-256 code hash is
\path{21ea83d80f459d1cac781bc8d2d3343330275bb1e13c3a0031f9331c246f32fb}.
It was naturally called in 23 matched task rows and recorded no
side-effect incident. Its called-subset mean paired outcome delta was
-0.0037, effectively neutral at this sample size. Lifecycle reflection
retained it with a route repair requirement because its applicability
attempts included frequent safe abstentions or failed selections and its
contribution was mixed.

\section{Tool 4: Message-Counterparty Contact Update}
\label{sec:appendix-message-counterparty-update}

\subsection{Purpose}\label{purpose-3}

This tool combines deterministic record selection with one contact
mutation. It receives visible message records, selects the unique latest
or oldest message, identifies the non-self counterparty, and applies
explicit update values to that person\textquotesingle s contact record.

\subsection{Execution Logic}\label{execution-logic-3}

\begin{enumerate}
\def\labelenumi{\arabic{enumi}.}
\tightlist
\item
  Abstain when records, updates, or the visible self-person identifier
  are missing.
\item
  Accept only \texttt{latest} or \texttt{oldest} as the selection mode.
\item
  Require every candidate message to contain a creation timestamp.
\item
  Select the unique timestamp extreme; abstain on a tie.
\item
  Compare sender and recipient identifiers with
  \texttt{self\_person\_id} to bind the non-self counterparty.
\item
  Remove update entries whose value is \texttt{None}.
\item
  Invoke \texttt{modify\_contact} exactly once for the selected
  counterparty.
\end{enumerate}

The generated tool receives the actor-visible records and self
identifier. It does not access hidden database state, expected benchmark
state, or a scenario-specific answer.

\subsection{Accepted Generated Code}\label{accepted-generated-code-3}

\begin{lstlisting}[style=sagepython,caption={Accepted message-counterparty contact-update native-action tool.},label={lst:sage-message-counterparty-update}]
def select_message_counterparty_for_contact_update(records: list, selection_mode: str, updates: dict, self_person_id: str):
    if not records or not updates or (not self_person_id):
        return {'status': 'abstain', 'confirmation': '', 'abstain_reason': 'Missing required inputs.', 'native_action': '', 'native_result': None}
    if selection_mode not in ['latest', 'oldest']:
        return {'status': 'abstain', 'confirmation': '', 'abstain_reason': 'Invalid selection mode.', 'native_action': '', 'native_result': None}
    if any((not isinstance(record, dict) or 'creation_timestamp' not in record for record in records)):
        return {'status': 'abstain', 'confirmation': '', 'abstain_reason': 'Invalid record structure.', 'native_action': '', 'native_result': None}
    winners = [record for record in records if 'creation_timestamp' in record]
    if not winners:
        return {'status': 'abstain', 'confirmation': '', 'abstain_reason': 'No valid records.', 'native_action': '', 'native_result': None}
    extreme_value = max(winners, key=lambda x: x['creation_timestamp']) if selection_mode == 'latest' else min(winners, key=lambda x: x['creation_timestamp'])
    winners = [record for record in records if record['creation_timestamp'] == extreme_value['creation_timestamp']]
    if len(winners) != 1:
        return {'status': 'abstain', 'confirmation': '', 'abstain_reason': 'Ambiguous selection.', 'native_action': '', 'native_result': None}
    selected_record = winners[0]
    person_id = selected_record['recipient_person_id'] if selected_record['recipient_person_id'] != self_person_id else selected_record['sender_person_id']
    if not person_id:
        return {'status': 'abstain', 'confirmation': '', 'abstain_reason': 'No non-self counterparty found.', 'native_action': '', 'native_result': None}
    native_result = modify_contact(person_id=person_id, **{k: v for k, v in updates.items() if v is not None})
    return {'status': 'success', 'confirmation': f'Updated contact for {person_id}.', 'abstain_reason': '', 'native_action': 'modify_contact', 'native_result': native_result}
\end{lstlisting}

\subsection{Validation and
Interpretation}\label{validation-and-interpretation-3}

The implementation passed one source example, two held-out examples,
four negative applicability examples, and runtime smoke validation. Its
SHA-256 code hash is
\path{fb2348049c1e52480ebd314cac9fa98bf0a2ed061edcee15b9773160d08f7e64}.
It was naturally called in 29 matched task rows, all 29 of which were
classified as helpful. The mean paired outcome delta was +0.5893, with
no runtime failure or side-effect incident. Lifecycle reflection
retained the tool.

\section{Tool 5: Reminder Recency Action}
\label{sec:appendix-recency-action-target}

\subsection{Purpose}\label{purpose-4}

This tool selects one reminder from visible records using deterministic
recency and equality constraints, then performs one explicitly requested
reminder action. It supports modification and removal while preventing
action under ambiguous selection.

\subsection{Execution Logic}\label{execution-logic-4}

\begin{enumerate}
\def\labelenumi{\arabic{enumi}.}
\tightlist
\item
  Default absent constraint and update mappings to empty dictionaries.
\item
  Accept only \texttt{modify\_reminder} or \texttt{remove\_reminder} as
  the action type.
\item
  Accept only \texttt{latest} or \texttt{oldest} as the selection mode.
\item
  Validate that all records contain the requested timestamp field.
\item
  Filter records by explicit equality constraints.
\item
  Select the unique minimum or maximum timestamp; abstain on no match or
  a tie.
\item
  Invoke \path{remove_reminder} once, or validate the update mapping
  and invoke \path{modify_reminder} once.
\end{enumerate}

The native action is selected by an explicit input covered by
validation. The tool does not infer whether modification or removal is
desired from hidden state.

\subsection{Accepted Generated Code}\label{accepted-generated-code-4}

\begin{lstlisting}[style=sagepython,caption={Accepted reminder-recency native-action tool.},label={lst:sage-recency-action-target}]
def select_action_target_by_recency(records: list=None, timestamp_key: str='', selection_mode: str='latest', action_type: str='', constraints: dict=None, updates: dict=None):
    if constraints is None:
        constraints = {}
    if updates is None:
        updates = {}
    if action_type not in ['modify_reminder', 'remove_reminder']:
        return {'status': 'abstain', 'confirmation': '', 'abstain_reason': 'Invalid action type.', 'native_action': '', 'native_result': None}
    if selection_mode not in ['latest', 'oldest']:
        return {'status': 'abstain', 'confirmation': '', 'abstain_reason': 'Invalid selection mode.', 'native_action': '', 'native_result': None}
    if not records or not all((isinstance(record, dict) and timestamp_key in record for record in records)):
        return {'status': 'abstain', 'confirmation': '', 'abstain_reason': 'No valid records found.', 'native_action': '', 'native_result': None}
    if not all((isinstance(record, dict) and timestamp_key in record for record in records)):
        return {'status': 'abstain', 'confirmation': '', 'abstain_reason': 'Invalid record structure.', 'native_action': '', 'native_result': None}
    filtered_records = [record for record in records if all((record.get(k) == v for k, v in constraints.items()))]
    if not filtered_records:
        return {'status': 'abstain', 'confirmation': '', 'abstain_reason': 'No records meet the constraints.', 'native_action': '', 'native_result': None}
    if selection_mode == 'latest':
        max_timestamp = max((record[timestamp_key] for record in filtered_records))
        winners = [record for record in filtered_records if record[timestamp_key] == max_timestamp]
    else:
        min_timestamp = min((record[timestamp_key] for record in filtered_records))
        winners = [record for record in filtered_records if record[timestamp_key] == min_timestamp]
    if len(winners) != 1:
        return {'status': 'abstain', 'confirmation': '', 'abstain_reason': 'Ambiguous selection of records.', 'native_action': '', 'native_result': None}
    selected_record = winners[0]
    if action_type == 'remove_reminder':
        reminder_id = selected_record['reminder_id']
        native_result = remove_reminder(reminder_id)
        return {'status': 'success', 'confirmation': f'Reminder {reminder_id} removed successfully.', 'abstain_reason': '', 'native_action': 'remove_reminder', 'native_result': native_result}
    elif action_type == 'modify_reminder':
        reminder_id = selected_record['reminder_id']
        if not updates:
            return {'status': 'abstain', 'confirmation': '', 'abstain_reason': 'No updates provided for modification.', 'native_action': '', 'native_result': None}
        reminder_timestamp = updates.get(timestamp_key, selected_record[timestamp_key])
        native_result = modify_reminder(reminder_id=reminder_id, reminder_timestamp=reminder_timestamp)
        return {'status': 'success', 'confirmation': f'Reminder {reminder_id} modified successfully.', 'abstain_reason': '', 'native_action': 'modify_reminder', 'native_result': native_result}
\end{lstlisting}

\subsection{Validation and
Interpretation}\label{validation-and-interpretation-4}

The implementation passed one source example, two held-out examples,
three negative applicability examples, and runtime smoke validation. Its
SHA-256 code hash is
\path{5dd847eb379de48f2556d0d3c75c2ce6801c24de7a42ed112f35b0f00c10d0cc}.
It was naturally called in 30 matched task rows and produced a mean
paired outcome delta of +0.6818. Nineteen rows were classified as
helpful and four as harmful. No runtime or side-effect incident was
recorded. Lifecycle reflection retained the tool with route repair
because its aggregate contribution was strongly positive but not
uniformly positive across routed contexts.

\section{Tool 6: Direct Contact Update}
\label{sec:appendix-contact-update-id}

\subsection{Purpose}\label{purpose-5}

This tool applies explicit updates to a contact that is already
identified by a visible person identifier. It validates that at least
one supported update field is present, normalizes phone-number
punctuation, requires international phone-number notation, and delegates
the mutation to \texttt{modify\_contact}.

\subsection{Execution Logic}\label{execution-logic-5}

\begin{enumerate}
\def\labelenumi{\arabic{enumi}.}
\tightlist
\item
  Abstain if \texttt{person\_id} is missing or neither a phone number
  nor name update is present.
\item
  Normalize spaces, hyphens, and parentheses from the supplied phone
  number.
\item
  Require a nonempty normalized number to begin with \texttt{+}.
\item
  Construct the native action arguments using only supplied update
  values.
\item
  Invoke \texttt{modify\_contact} exactly once.
\end{enumerate}

The tool does not search for or infer the contact identifier. The
identifier must already be visible to the actor through the ordinary
task interaction.

\subsection{Accepted Generated Code}\label{accepted-generated-code-5}

\begin{lstlisting}[style=sagepython,caption={Accepted direct-contact-update native-action tool.},label={lst:sage-contact-update-id}]
def plan_contact_update_from_id(person_id: str=None, phone_number: str=None, name: str='', relationship: str='', user_request: str=''):
    if not person_id or (not phone_number and (not name)):
        return {'status': 'abstain', 'native_action': '', 'native_result': None, 'abstain_reason': 'Missing person_id or no update fields provided.', 'confirmation': ''}
    normalized_phone_number = phone_number.replace(' ', '').replace('-', '').replace('(', '').replace(')', '')
    if not normalized_phone_number.startswith('+') and normalized_phone_number:
        return {'status': 'abstain', 'native_action': '', 'native_result': None, 'abstain_reason': 'Phone number must start with a plus sign.', 'confirmation': ''}
    kwargs = {'person_id': person_id}
    if normalized_phone_number:
        kwargs['phone_number'] = normalized_phone_number
    if name:
        kwargs['name'] = name
    native_result = modify_contact(**kwargs)
    return {'status': 'success', 'native_action': 'modify_contact', 'native_result': native_result, 'abstain_reason': '', 'confirmation': f'Updated contact with ID {person_id} successfully.'}
\end{lstlisting}

\subsection{Validation and
Interpretation}\label{validation-and-interpretation-5}

The implementation passed one source example, one held-out example, one
negative applicability example, and runtime smoke validation. Its
SHA-256 code hash is
\path{8271ed9488e1d30dbb1290cd3ca0e8c74b60e7c0ba636efef099fc908c1a04e7}.
It was naturally called in seven matched task rows with no runtime
failure or side-effect incident. Its mean paired outcome delta was
-0.0476, and lifecycle reflection classified it as requiring routing
repair. The accepted code demonstrates valid bounded native delegation,
but its observed routing did not provide positive contribution evidence
in this run.

\section{Methodological
Interpretation}\label{methodological-interpretation}

These implementations demonstrate that SAGE\textquotesingle s generated
tools were not limited to advisory text or argument suggestions. SAGE
generated executable tools that combined deterministic visible-evidence
processing with preserved environment actions. The state changes
remained attributable to native ToolSandbox calls, and every generated
implementation retained an explicit abstention path for incomplete,
unsupported, or ambiguous inputs.

The evidence also shows why structural validity and empirical utility
must be reported separately. All six tools satisfied the native-action
admission contract, but only three produced clearly positive
called-subset outcome deltas in this run.
\path{prepare_reminder_creation_args},
\path{select_message_counterparty_for_contact_update}, and
\path{select_action_target_by_recency} provide the strongest
evidence that a generated native-action tool can improve downstream task
completion. \path{apply_single_device_state_action} was
approximately neutral, while \texttt{prepare\_add\_contact\_args} and
\path{plan_contact_update_from_id} require improved routing before
they should be promoted as performance-enhancing tools.

Accordingly, the defensible conclusion is not that every autonomously
generated native-action tool improves performance. The evidence supports
the narrower and stronger claim that SAGE can autonomously generate,
validate, store, naturally call, and reuse bounded tools that execute
preserved native actions, and that selected generated tools produce
substantial positive outcome differences on matched tasks.

\section{Reproducibility and Evidence
Boundary}\label{reproducibility-and-evidence-boundary}

The code listings above are copied from the immutable accepted-tool
records in the registry manifest. The associated code hashes permit
exact verification. The lifecycle metrics are copied from the
corresponding post-run contribution record. No benchmark expected-answer
strings, hidden task labels, scenario identifiers, or synthetic
code-generated final responses are present in the tool implementations.

The formal run used the sealed 1,032-task manifest with SHA-256 digest
\path{21877bd3524258b80f74207c66ed3640b6db629d13b4a2fb4d817e35d0390bec}.
The baseline used 1,032 strict control-cache records; the SAGE arm was
fresh and began with an empty registry. The run completed with zero
runtime exceptions. The full run produced a baseline task-completion
outcome of 0.4573 and a SAGE outcome of 0.7900, corresponding to an
absolute outcome difference of +0.3326 and a relative outcome lift of
+72.73 percent.
